\documentclass[oneside,a4paper,14pt]{extarticle}
\usepackage[T1,T2A]{fontenc}
\usepackage[a4paper,top=20mm,bottom=20mm,left=30mm,right=10mm]{geometry}
\usepackage{amsmath, amssymb}
\usepackage[utf8]{inputenc}
\usepackage[russian]{babel}
\usepackage{textcomp}
\usepackage{indentfirst}
\usepackage{graphicx}
\usepackage{float}
\usepackage{mwe}
\usepackage{wrapfig}
\usepackage{caption}
\usepackage{tocloft}
\renewcommand{\cftsecleader}{\cftdotfill{\cftdotsep}}
\usepackage{longtable}
\usepackage{amsfonts}
\usepackage{amsthm}
\usepackage[all]{xy}
\usepackage[breaklinks]{hyperref}
\usepackage{titlesec}
\usepackage{fancyvrb}
\usepackage{paracol}
\usepackage{listings}
\usepackage{xcolor}
\usepackage{listings}

\begin{document}

\newpage\thispagestyle{empty}
\begin{center}
    МИНИСТЕРСТВО НАУКИ И ВЫСШЕГО ОБРАЗОВАНИЯ\\
    РОССИЙСКОЙ ФЕДЕРАЦИИ\\
    ФЕДЕРАЛЬНОЕ ГОСУДАРСТВЕННОЕ БЮДЖЕТНОЕ\\
    ОБРАЗОВАТЕЛЬНОЕ УЧРЕЖДЕНИЕ ВЫСШЕГО ОБРАЗОВАНИЯ\\
    «ВЯТСКИЙ ГОСУДАРСТВЕННЫЙ УНИВЕРСИТЕТ»\\
    Институт математики и информационных систем\\
    Факультет автоматики и вычислительной техники\\
    Кафедра электронных вычислительных машин
\end{center}

\vspace{20mm}
\begin{center}

    Отчёт по лабораторной работе №1\\
    по дисциплине\\
    «Вычислительная математика и методы оптимизации»\\
\end{center}
\vfill

Выполнил студент гр. ИВТб-2301-05-00 \hfill
\makebox[8cm]{\hrulefill /Ковальский И.И./}

Проверил преподаватель\hfill
\makebox[8cm]{\hrulefill /Коржавина А.С./}

\vfill
\begin{center}
    Киров\\
    2026
\end{center}

\newpage
\setcounter{page}{1}
\tableofcontents
\newpage

\section{Задание}
\begin{enumerate}
    \item Изучить и реализовать заданный метод. Язык программирования выбрать самостоятельно.
    \item Продемонстрировать работу приложения на примерах. При демонстрации работы приложения выводить каждую итерацию метода (симплекс-таблицы) отдельно, также выводить решение (ответ). Предусмотреть все возможные варианты решения задачи, включая отсутствие решений и т.п.
\end{enumerate}
\textbf{Заданный метод} --- двойственный симплекс-метод.

\section{Листинг программы}

\lstinputlisting[language=Python]{1lab.py}

\newpage

\section{Результат работы}
\subsection{Пример 1: Целевая функция не ограничена}
\textbf{Задача:}
$F = 2x_1 + 3x_2 \to \max$ \\
$-2x_1 - 3x_2 \le -6$ \\
$-x_1 - x_2 \le -1$ \\
$-x_1 + 2x_2 \le 2$ \\

\textbf{Вывод программы:}
\begin{verbatim}
=================================================================
  ДВОЙСТВЕННЫЙ СИМПЛЕКС-МЕТОД
=================================================================
  Начальная таблица
Базис          x1       x2       s1       s2       s3        b
-----------------------------------------------------------------
s1         -2.000   -3.000    1.000    0.000    0.000   -6.000
s2         -1.000   -1.000    0.000    1.000    0.000   -1.000
s3         -1.000    2.000    0.000    0.000    1.000    2.000
F          -2.000   -3.000    0.000    0.000    0.000    0.000
-----------------------------------------------------------------
... (промежуточные таблицы) ...
  ЦЕЛЕВАЯ ФУНКЦИЯ НЕ ОГРАНИЧЕНА
  — задача не имеет конечного оптимума.
\end{verbatim}

\subsection{Пример 2: Нет допустимых решений}
\textbf{Задача:}
$F = x_1 + x_2 \to \max$ \\
$-x_1 \le -3$ (то есть $x_1 \ge 3$) \\
$x_1 \le 1$ \\

\textbf{Вывод программы:}
\begin{verbatim}
=================================================================
  ДВОЙСТВЕННЫЙ СИМПЛЕКС-МЕТОД
=================================================================
  Начальная таблица
Базис          x1       x2       s1       s2        b
-----------------------------------------------------------------
s1         -1.000    0.000    1.000    0.000   -3.000
s2          1.000    0.000    0.000    1.000    1.000
F          -1.000   -1.000    0.000    0.000    0.000
-----------------------------------------------------------------
  ЗАДАЧА НЕ ИМЕЕТ РЕШЕНИЙ:
  система ограничений несовместна.
\end{verbatim}

\section{Вывод}
В ходе данной лабораторной работы был изучен и реализован двойственный симплекс-метод для решения задач линейного программирования. Разработанное приложение демонстрирует пошаговое выполнение алгоритма с выводом симплекс-таблиц на каждой итерации.

Программа была протестирована на различных примерах, что позволило убедиться в ее корректной работе в следующих ситуациях:
\begin{enumerate}
    \item Нахождение оптимального решения.
    \item Определение несовместности системы ограничений.
    \item Выявление неограниченности целевой функции.
\end{enumerate}
Поставленная задача выполнена в полном объеме.

\end{document}